\subsection{Transition to Product Development}
\label{subsec:ueberfuehrung-produktentwicklung}

After generation, the independent product project assumes responsibility for
domain models, business rules, the user interface, integrations, data, and
operations. Authentication, authorization, scaling, monitoring, backups, data
protection, and signing also remain product tasks. The baseline structures
this work, but does not solve these product-specific tasks.

The master repository and the product have separate lifecycles. There is no
permanent update or synchronization mechanism. Changes to the master
repository must be transferred to products deliberately; product changes
should flow back into the baseline only after a generalization and
verification review.

Template acceptance and product acceptance remain separate. The
commit-specific acceptance provides evidence for generation, installation,
tests, and selected baseline builds, but not for deployment, signing, or
functional and security requirements
\parencite{kleiveist2026acceptance}. The product requires its own tests and
evidence for these areas.
They must cover the specific customer requirements, integrations, and
operating conditions.
\beginnerinput{workspace/statement/800_einsatz_transfer/860}
